\begin{statementbox}
\textbf{\textcolor{CStatement}{核心结论：用同一个角表示两个坐标}}
设椭圆中心为原点，方程为
\[
\frac{x^2}{a^2}+\frac{y^2}{b^2}=1,\qquad a>b>0.
\]
$a,b$ 分别为横向、纵向半轴长。椭圆上的点恰好可表示为
\[
\boxed{P(a\cos\theta,b\sin\theta),\qquad0\le\theta<2\pi.}
\]
每个参数给出的点都在椭圆上，椭圆上每个点也都能取到。

\begin{center}
\begin{tikzpicture}[x=1cm,y=1cm,font=\small,>=Stealth]
\draw[rounded corners=3pt,gray!45] (-6.3,-4.05) rectangle (6.3,4.05);
\node[font=\small\bfseries] at (0,3.5999999999999996) {同一横坐标：圆上的点压到椭圆上};
\draw[->,gray] (-3.5,0)--(3.7,0) node[right] {$x$};\draw[->,gray] (0,-3.15)--(0,3.15) node[right] {$y$};\node[below left] at (0,0) {$O$};\draw[orange!80!black,dashed,thick] (0,0) circle (3);\draw[CStatement,very thick] (0,0) ellipse (3 and 1.8);
\coordinate (Q) at ({3*cos(55)},{3*sin(55)});
\coordinate (P) at ({3*cos(55)},{1.8*sin(55)});
\draw[orange!80!black,thick] (0,0)--(Q);
\draw[->,CProof,thick] (Q)--(P);
\draw[densely dotted] (P)--({3*cos(55)},0);
\draw[->,orange!80!black] (.7,0) arc (0:55:.7);
\node at (.95,.45) {$\theta$};
\fill[orange!80!black] (Q) circle (2pt) node[above right] {$Q(a\cos\theta,a\sin\theta)$};
\fill[CStatement] (P) circle (2pt) node[above right] {$P$};\node[anchor=west,align=left] at (3.25,-1.4) {$P$ 的坐标\\$(a\cos\theta,b\sin\theta)$};
\node[below right] at (3.08,-.08) {$a$};\node[above left] at (0,1.8) {$b$};
\node[anchor=east,align=right] at (-3.45,1.6) {虚线：辅助圆\\半径 $a$};
\node[anchor=east,align=right,text=CStatement] at (-3.45,-1.4) {实线：椭圆\\纵坐标乘 $b/a$};

\end{tikzpicture}
\end{center}


\textbf{看图记来路：}保持圆上点 $Q$ 的横坐标，将纵坐标乘 $b/a$，得到椭圆上的 $P$；$\theta$ 画在 $OQ$ 与 $x$ 轴正方向之间。

\tcbbreak
\textbf{\textcolor{CStatement}{方向与范围}}
也可让 $\theta\in\mathbb R$，但每隔 $2\pi$ 重复同一点。若长轴竖直，
\[
\frac{x^2}{b^2}+\frac{y^2}{a^2}=1
\quad\Longrightarrow\quad P(b\cos\theta,a\sin\theta).
\]
横坐标前的系数对应横向半轴，纵坐标前的系数对应纵向半轴。

\medskip
\textbf{\textcolor{CStatement}{中心到椭圆的距离}}
对主式中的椭圆，$O=(0,0)$ 时有
\[
\boxed{b\le OP\le a.}
\]
最短在 $(0,\pm b)$，最长在 $(\pm a,0)$；同心圆与椭圆的接触图见理解说明。

\medskip
\textbf{\textcolor{CStatement}{提醒}}
$\theta$ 来自辅助圆，一般不等于 $OP$ 与 $x$ 轴正方向的夹角。若只取某段椭圆，必须同时限制 $\theta$。
\end{statementbox}
